\problemname{Chokoladeæsken}
Alex elsker chokolade.
Han har derfor altid en æske i spisekammeret, som han sommetider smugspiser fra.
Når æsken er tom, køber han en ny og lader som ingenting.
Alex' kæreste Kim aner dog uråd -- hvorfor bliver den der æske egentlig aldrig tom? -- og er begyndt at holde øje med antallet af de resterende chokoladestykker.

Skriv et program, der på grundlag af Kims observationer udregner det mindste antal nye æsker, som Alex må have købt i løbet af perioden.

\section*{Indlæsning}
På første linje står et heltal $N$ ($1 \le N \le 100$), antallet af observationer.

Derefter følger en linje med $N$ heltal $a_1, \ldots, a_N$, antallet af chokoladestykker i æsken ved hver observation, i den rækkefølge de foretages.
Der gælder $1 \le a_i \le 100$ for alle $i$.

\section*{Udskrift}
Skriv en enkelt linje indeholdende et heltal: det mindste antal nye chokoladeæsker, som Alex beviseligt må have købt i løbet af perioden.
